% -*- coding: utf-8 -*-
% !TEX program = xelatex
\documentclass[plain,noanswer,medmath]{../../template/jnuexam}
\usepackage{../../template/kysx}

\begin{document}


\examtitle{2004}{2}


\loadproblems{1}{2004P1}%载入试卷一


\makepart{填空题}{1～6小题，每小题4分，共24分}


\begin{problem}
设$f(x) = \lim_{n \to \infty} \frac{(n - 1)x}{nx^2 + 1}$，则$f(x)$的间断点为$x =$ \fillin{}．
\end{problem}


\begin{problem}
设函数$y(x)$由参数方程$\begin{cases}
  x = t^3 + 3t + 1 \\
  y = t^3 - 3t + 1
\end{cases}$确定，则曲线$y = y(x)$向上凸的$x$取值范围为 \fillin{}．
\end{problem}


\begin{problem}
$\int_{1}^{+\infty}\!\!\frac{\dx}{x\sqrt{x^2 - 1}}=$ \fillin{}．
\end{problem}


\begin{problem}
设函数$z = z(x,y)$由方程$z = \e^{2x - 3z} + 2y$确定，则$3\frac{\pdz}{\pdx} + \frac{\pdz}{\pdy} =$ \fillin{}．
\end{problem}


\begin{problem}
微分方程$(y + x^3)\dx - 2x\dy = 0$满足$\left. y \right|_{x = 1} = \frac{6}{5}$的特解为 \fillin{}．
\end{problem}


\useproblem{1}{1}{5}%【同试卷一第一(5)题】


\makepart{选择题}{7～14小题，每小题4分，共32分}


\useproblem{1}{2}{7}%【同试卷一第二(7)题】


\begin{problem}
设$f(x) = |x(1 - x)|$，则\pickout{}
\begin{abcd}
\item $x = 0$是$f(x)$的极值点，但$(0,0)$不是曲线$y = f(x)$的拐点．
\item $x = 0$不是$f(x)$的极值点，但$(0,0)$是曲线$y = f(x)$的拐点．
\item $x = 0$是$f(x)$的极值点，且$(0,0)$是曲线$y = f(x)$的拐点．
\item $x = 0$不是$f(x)$的极值点，$(0,0)$也不是曲线$y = f(x)$的拐点．
\end{abcd}
\end{problem}


\begin{problem}
$\lim_{n \to \infty} \ln \sqrt[n]{\left(1 + \frac{1}{n} \right)^2\left(1 + \frac{2}{n} \right)^2
 \cdots \left(1 + \frac{n}{n} \right)^2}$等于\pickout{}
\begin{abcd}
\item $\int_1^2 \ln^2 x\dx$．
\item $2\int_1^2 \ln x\dx$．
\item $2\int_1^2 \ln(1 + x)\dx$．
\item $\int_1^2 \ln^2(1 + x)\dx$．
\end{abcd}
\end{problem}


\useproblem{1}{2}{8}%【同试卷一第二(8)题】


\begin{problem}
微分方程$y'' + y = x^2 + 1 + \sin x$的特解形式可设为\pickout{}
\begin{abcd}
\item $y^* = ax^2 + bx + c + x(A\sin x + B\cos x)$．
\item $y^* = x(ax^2 + bx + c + A\sin x + B\cos x)$．
\item $y^* = ax^2 + bx + c + A\sin x$．
\item $y^* = ax^2 + bx + c + A\cos x$．
\end{abcd}
\end{problem}


\begin{problem}
设函数$f(u)$连续，区域$D = \{(x,y)|x^2 + y^2 \le 2y\}$，则$\iint_D f(xy)\dx\dy$等于\pickout{}
\begin{abcd}
\item $\int_{-1}^1 \dx \int_{- \sqrt{1 - x^2}}^{\sqrt{1 - x^2}} f(xy)\dy$.
\item $2\int_0^2 \dy \int_0^{\sqrt{2y - y^2}} f(xy)\dx$.
\item $\int_0^{\pi}\d\theta \int_0^{2\sin \theta} f(r^2\sin \theta \cos \theta)\dr$.
\item $\int_0^{\pi}\d\theta \int_0^{2\sin \theta} f(r^2\sin \theta \cos \theta)r\dr$.
\end{abcd}
\end{problem}


\useproblem{1}{2}{11}%【同试卷一第二(11)题】


\useproblem{1}{2}{12}%【同试卷一第二(12)题】


\makepart{解答题}{15～23小题，共94分}


\begin{problem}[points=10]
求极限$\lim_{x \to 0} \frac{1}{x^3}\left[\left(\frac{2 + \cos x}{3} \right)^x - 1 \right]$．
\end{problem}


\begin{problem}[points=10]
设函数$f(x)$在$(-\infty , +\infty)$上有定义，在区间$[0,2]$上，$f(x) = x(x^2 - 4)$，
若对任意的$x$都满足$f(x) = kf(x + 2)$，其中$k$为常数．
\begin{enumerate}
\item 写出$f(x)$在$[- 2,0]$上的表达式；
\item 问$k$为何值时，$f(x)$在$x = 0$处可导．
\end{enumerate}
\end{problem}


\begin{problem}[points=11]
设$f(x) = \int_x^{x + \frac{\pi}{2}} |\sin t|\dt$．\par
\begin{enumerate*}
\item 证明$f(x)$是以$\pi$为周期的周期函数；
\item 求$f(x)$的值域．
\end{enumerate*}
\end{problem}


\begin{problem}[points=12]
曲线$y = \frac{\e^x + \e^{- x}}{2}$与直线$x = 0$, $x = t$（$t > 0$）及$y = 0$围成一曲边梯形．
该曲边梯形绕$x$轴旋转一周得一旋转体，其体积为$V(t)$，侧面积为$S(t)$，在$x = t$处的底面积为$F(t)$．\par
\begin{enumerate*}
\item 求$\frac{S(t)}{V(t)}$的值；
\item 计算极限$\lim_{t \to +\infty} \frac{S(t)}{F(t)}$．
\end{enumerate*}
\end{problem}


\useproblem[points=12]{1}{3}{15}%【同试卷一第三(15)题】


\useproblem[points=11]{1}{3}{16}%【同试卷一第三(16)题】


\begin{problem}[points=10]
设$z = f(x^2 - y^2,\e^{xy})$，其中$f$具有连续二阶偏导数，
求$\frac{\pdz}{\pdx}$，$\frac{\pdz}{\pdy}$，$\frac{\pd^2 z}{\pdx\pd y}$．
\end{problem}


\begin{problem}[points=9]%【类似试卷一第三(20题)】
设有齐次线性方程组
$$\left\{\begin{array}{*{20}{l}}
  \begin{array}{*{20}{l}}
  (1 + a) x_1 + x_2 + x_3 + x_4 = 0, \\
  2x_1 + (2 + a) x_2 + 2x_3 + 2x_4 = 0, \\
  3x_1 + 3x_2 + (3 + a) x_3 + 3x_4 = 0, \\
  4x_1 + 4x_2 + 4x_3 + (4 + a) x_4 = 0.
\end{array}
\end{array} \right.$$
试问$a$取何值时，该方程组有非零解，并求出其通解．
\end{problem}


\useproblem[points=9]{1}{3}{21}%【同试卷一第三(21)题】


\end{document}
